\section{Conceptos previos}
\textcolor{teal}{¿Introduzco un capítulo/sección 'Conceptos previos' que haga una introducción a los conceptos sobre los que se cimenta este trabajo? P.e.: Inteligencia Artificial, LLMs y otro tipo de conceptos de imprescindible conocimiento para la comprensión del trabajo.}

\textcolor{teal}{
    Qué incluir:
    \begin{itemize}
        \item Lenguaje natural y evaluación
        \item Qué es el lenguaje natural (texto y habla)
    \end{itemize}
    Qué implica “evaluar lenguaje”:
    \begin{itemize}
        \item forma (gramática, ortografía)
        \item contenido (semántica, adecuación)
        \item pragmática (intención, contexto)
        \item Modelos de lenguaje (LLMs)
        \item Definición de LLM
    \end{itemize} 
    Capacidades:
    \begin{itemize}
        \item generación
        \item comprensión
        \item evaluación (emergente)
        \item Idea clave:
            Los LLMs no fueron diseñados explícitamente para evaluar, pero muestran capacidad para hacerlo
    \end{itemize}
    ASR (Automatic Speech Recognition)
    \begin{itemize}
        \item Qué es
        \item Pipeline básico:
        \item audio → texto
    \end{itemize}
    Problemas:
    \begin{itemize}
        \item ruido
        \item acentos
        \item niños (importante para tu contexto)
    \end{itemize}
    Cierre del apartado:
    Introduce la idea de que evaluar lenguaje automático implica múltiples niveles → necesidad de enfoques complesjos
}
